% !TEX program = xelatex
\documentclass[11pt,a4paper]{article}

\usepackage[UTF8,fontset=none]{ctex}
\setCJKmainfont[BoldFont=Heiti SC,ItalicFont=Songti SC]{Songti SC}
\setCJKsansfont{Heiti SC}
\setCJKmonofont{Heiti SC}

\usepackage{amsmath,amssymb,amsfonts,mathtools}
\usepackage{booktabs,array,tabularx}
\usepackage[margin=1.9cm]{geometry}
\usepackage[dvipsnames]{xcolor}
\usepackage[most]{tcolorbox}
\usepackage[hidelinks]{hyperref}
\usepackage{enumitem}
\usepackage{fancyhdr}

\setlength{\parindent}{0em}
\setlength{\parskip}{0.45em}
\setlist[itemize]{leftmargin=*,topsep=2pt,itemsep=1.5pt}
\setlist[enumerate]{leftmargin=*,topsep=2pt,itemsep=1.5pt}
\allowdisplaybreaks

\pagestyle{fancy}
\fancyhf{}
\lhead{\small Statistical Modelling and Inference}
\rhead{\small 2026 第二套模拟样卷}
\cfoot{\thepage}

\hypersetup{
  pdftitle={SMI Final Sample Examination 2 with Detailed Solutions 2026},
  pdfauthor={Prepared from the confirmed exam scope, lecture examples, assignments, and review notes}
}

\newcommand{\E}{\operatorname{E}}
\newcommand{\Var}{\operatorname{Var}}
\newcommand{\Cov}{\operatorname{Cov}}
\newcommand{\Bias}{\operatorname{Bias}}
\newcommand{\MSE}{\operatorname{MSE}}
\newcommand{\SE}{\operatorname{SE}}
\newcommand{\iid}{\overset{\mathrm{iid}}{\sim}}
\newcommand{\by}{\mathbf{y}}
\newcommand{\bX}{\mathbf{X}}
\newcommand{\bH}{\mathbf{H}}
\newcommand{\bI}{\mathbf{I}}
\newcommand{\be}{\mathbf{e}}
\newcommand{\bbeta}{\boldsymbol{\beta}}
\newcommand{\bvarepsilon}{\boldsymbol{\varepsilon}}
\newcommand{\pts}[1]{\hfill\textbf{[#1 分]}}

\newtcolorbox{instructionbox}[1][]{
  colback=gray!5!white,colframe=gray!65!black,
  fonttitle=\bfseries,title=#1,breakable
}
\newtcolorbox{questionbox}[1][]{
  colback=blue!3!white,colframe=blue!60!black,
  fonttitle=\bfseries,title=#1,breakable
}
\newtcolorbox{solutionbox}[1][]{
  colback=green!3!white,colframe=green!55!black,
  fonttitle=\bfseries,title=#1,breakable
}
\newtcolorbox{checkbox}[1][]{
  colback=orange!5!white,colframe=orange!70!black,
  fonttitle=\bfseries,title=#1,breakable
}
\newtcolorbox{warnbox}[1][]{
  colback=red!4!white,colframe=red!60!black,
  fonttitle=\bfseries,title=#1,breakable
}

\title{\textbf{Statistical Modelling and Inference}\\[0.3em]
2026 期末模拟样卷（二）与详细解析}
\author{依据最新考试范围、课件例题、Assignment 1/2 与复习笔记编制}
\date{June 2026}

\begin{document}
\maketitle

\begin{instructionbox}[考试说明]
\begin{itemize}
  \item 建议时长：120 分钟；满分：100 分；共五道大题，每题 20 分。
  \item 必须展示主要计算与推导过程。只有最终答案、没有方法时可能无法获得完整分数。
  \item 假设检验统一使用\textbf{临界值法}，不要求计算 P-value。
  \item 本卷覆盖：Bias/MSE、单均值与比例推断、独立/配对两均值、简单回归手算、多元回归理论证明、Normal/Binomial/Poisson/Exponential MLE。
  \item 本卷不涉及：ANOVA/ANCOVA、R 编程、P-value、Power、polynomial regression、nested-model comparison。
\end{itemize}
\end{instructionbox}

\begin{instructionbox}[本卷可直接使用的临界值]
\[
  z_{0.05}=1.645,
  \qquad
  z_{0.025}=1.960.
\]
\[
\begin{aligned}
  t_{0.05,24}&=1.711,
  &t_{0.025,24}&=2.064,\\
  t_{0.05,15}&=1.753,
  &t_{0.025,15}&=2.131,\\
  t_{0.025,8}&=2.306.
\end{aligned}
\]
\end{instructionbox}

\tableofcontents
\newpage

% ============================================================
\section{模拟试卷}
% ============================================================

\begin{questionbox}[第 1 题：Bias、MSE 与无偏修正（20 分）]
设 $T$ 是参数 $\theta$ 的估计量，并且
\[
  \E(T)=2\theta+1,
  \qquad
  \Var(T)=\frac9n.
\]

\begin{enumerate}
  \item 从定义出发证明
  \[
    \MSE(\hat\theta)
    =
    \Var(\hat\theta)+\Bias(\hat\theta)^2.
  \]
  \pts{5}

  \item 求 $T$ 作为 $\theta$ 的估计量时的 bias 与 MSE。 \pts{4}

  \item 构造形如 $U=aT+b$ 的 $\theta$ 的无偏估计量。证明其无偏，并求
  $\Var(U)$ 与 $\MSE(U)$。 \pts{5}

  \item 另有一个与 $U$ 独立的无偏估计量 $S$，满足
  \[
    \Var(S)=\frac3n.
  \]
  对
  \[
    W_c=cU+(1-c)S
  \]
  求使 MSE 最小的 $c$，并求最小 MSE。 \pts{6}
\end{enumerate}
\end{questionbox}

\begin{questionbox}[第 2 题：单总体均值与总体比例（20 分）]
\textbf{第一部分：总体均值。}

某正态总体的均值 $\mu$ 与方差均未知。随机抽取 $n=25$ 个观测，得到
\[
  \bar x=41.2,
  \qquad
  s=5.0.
\]

\begin{enumerate}
  \item 构造 $\mu$ 的 95\% 置信区间，并作语境解释。 \pts{5}

  \item 使用临界值法检验
  \[
    H_0:\mu\leq40
    \qquad\text{vs.}\qquad
    H_a:\mu>40.
  \]
  写出统计量、参考分布、拒绝域、决定和结论。 \pts{7}
\end{enumerate}

\textbf{第二部分：总体比例。}

随机调查 $400$ 人，其中 $116$ 人具有某项特征。令 $p$ 表示总体中具有该特征的比例。

\begin{enumerate}[resume]
  \item 构造 $p$ 的 95\% Wald 置信区间。 \pts{4}

  \item 使用临界值法检验
  \[
    H_0:p=0.25
    \qquad\text{vs.}\qquad
    H_a:p\ne0.25.
  \]
  明确说明为什么检验标准误和置信区间标准误中的比例不同。 \pts{4}
\end{enumerate}
\end{questionbox}

\begin{questionbox}[第 3 题：配对样本与两个独立大样本（20 分）]
\textbf{第一部分：配对样本。}

同一批 $16$ 名学生在教学改革前后各参加一次测试。令
\[
  D_i=(\text{改革后成绩})-(\text{改革前成绩}).
\]
差值摘要为
\[
  \bar d=2.4,
  \qquad
  s_d=3.0,
  \qquad
  n=16.
\]

\begin{enumerate}
  \item 解释为什么本题应使用 paired-sample 方法，而不是把两次成绩当成两个独立样本。 \pts{3}

  \item 使用临界值法检验教学改革是否提高平均成绩：
  \[
    H_0:\mu_D\leq0
    \qquad\text{vs.}\qquad
    H_a:\mu_D>0.
  \]
  \pts{4}

  \item 构造 $\mu_D$ 的 95\% 置信区间并解释。 \pts{3}
\end{enumerate}

\textbf{第二部分：两个独立大样本。}

两个独立总体的样本摘要为
\[
\begin{array}{c|ccc}
 & n & \bar x & s\\
\hline
\text{总体 1} & 100 & 12.8 & 4.0\\
\text{总体 2} & 81 & 11.5 & 3.6
\end{array}
\]

\begin{enumerate}[resume]
  \item 构造 $\mu_1-\mu_2$ 的 95\% 置信区间。 \pts{5}

  \item 使用临界值法检验
  \[
    H_0:\mu_1-\mu_2=0
    \qquad\text{vs.}\qquad
    H_a:\mu_1-\mu_2\ne0.
  \]
  \pts{5}
\end{enumerate}
\end{questionbox}

\begin{questionbox}[第 4 题：简单回归计算与多元回归证明（20 分）]
\textbf{第一部分：简单线性回归。}

对 $n=10$ 个观测拟合
\[
  Y_i=\beta_0+\beta_1x_i+\varepsilon_i.
\]
已知
\[
  \bar x=12,\qquad
  \bar y=30,\qquad
  S_{xx}=80,\qquad
  S_{xy}=200,\qquad
  S_{yy}=620.
\]

\begin{enumerate}
  \item 求 $\hat\beta_0,\hat\beta_1$，并写出拟合直线。再求
  $SSE$、$s_e^2$ 与 $R^2$。 \pts{6}

  \item 使用临界值法检验
  \[
    H_0:\beta_1=0
    \qquad\text{vs.}\qquad
    H_a:\beta_1\ne0.
  \]
  \pts{3}

  \item 当 $x_0=14$ 时，分别构造平均响应
  $\E(Y\mid x_0)$ 的 95\% CI 和一个新个体响应的 95\% prediction interval，并解释二者宽度不同的原因。 \pts{5}
\end{enumerate}

\textbf{第二部分：多元线性回归理论。}

考虑
\[
  \by=\bX\bbeta+\bvarepsilon,
  \qquad
  \E(\bvarepsilon)=\mathbf0,
  \qquad
  \Var(\bvarepsilon)=\sigma^2\bI,
\]
其中 $\bX$ 满列秩。

\begin{enumerate}[resume]
  \item 从最小化
  \[
    Q(\mathbf b)=(\by-\bX\mathbf b)^\mathsf{T}(\by-\bX\mathbf b)
  \]
  推导
  \[
    \hat{\bbeta}=(\bX^\mathsf{T}\bX)^{-1}\bX^\mathsf{T}\by,
  \]
  并说明为什么残差满足 $\bX^\mathsf{T}\be=\mathbf0$。 \pts{3}

  \item 证明
  \[
    \E(\hat{\bbeta})=\bbeta,
    \qquad
    \Var(\hat{\bbeta})
    =
    \sigma^2(\bX^\mathsf{T}\bX)^{-1}.
  \]
  \pts{3}
\end{enumerate}
\end{questionbox}

\begin{questionbox}[第 5 题：四类分布的 MLE 与 Fisher 信息（20 分）]
以下各小题相互独立。

\begin{enumerate}
  \item 设
  \[
    Y_1,\ldots,Y_8\iid N(\mu,\sigma^2),
  \]
  且
  \[
    \bar y=5,
    \qquad
    \sum_{i=1}^8(y_i-\bar y)^2=24.
  \]
  写出 log-likelihood 的主要形式，求 $\mu$ 与 $\sigma^2$ 的 MLE，并解释方差 MLE 为什么使用分母 $n$。 \pts{5}

  \item 设 $Y\sim\operatorname{Binomial}(20,p)$，观测到 $y=0$。求 $p$ 的 MLE。说明为什么必须检查参数空间边界，以及普通 Wald 标准误在本例中为什么不可靠。 \pts{5}

  \item 设
  \[
    Y_1,\ldots,Y_{16}\iid\operatorname{Poisson}(\lambda),
    \qquad
    \sum_{i=1}^{16}y_i=64.
  \]
  求 $\hat\lambda_{\mathrm{MLE}}$、样本 Fisher 信息和 plug-in 标准误，并构造 95\% Wald 置信区间。 \pts{5}

  \item 设
  \[
    Y_1,\ldots,Y_{16}\iid\operatorname{Exponential}(\lambda),
  \]
  其中 $\lambda$ 是 rate，且 $\sum y_i=40$。求 MLE 和 Fisher 信息，并使用 plug-in Wald 临界值法检验
  \[
    H_0:\lambda=0.6
    \qquad\text{vs.}\qquad
    H_a:\lambda\ne0.6.
  \]
  \pts{5}
\end{enumerate}
\end{questionbox}

\newpage
% ============================================================
\section{详细解析}
% ============================================================

\begin{checkbox}[本卷主要考查的答题动作]
\begin{enumerate}
  \item Bias/MSE 题先计算期望，再决定是否可以把 MSE 简化为方差。
  \item 区间与检验题必须区分：区间标准误使用样本估计，检验标准误通常按 $H_0$ 计算。
  \item 两均值题先判断 independent 还是 paired，再决定标准误。
  \item 回归计算题先算 $S_{xx},S_{xy}$ 路线；矩阵证明题先写正规方程。
  \item MLE 题必须写参数空间，并检查内部临界点、二阶导或边界。
\end{enumerate}
\end{checkbox}

\begin{solutionbox}[第 1 题解析：Bias、MSE 与无偏修正]
\textbf{（1）证明 bias--variance decomposition。}

令
\[
  b=\Bias(\hat\theta)=\E(\hat\theta)-\theta.
\]
在估计误差中加减 $\E(\hat\theta)$：
\[
  \hat\theta-\theta
  =
  \{\hat\theta-\E(\hat\theta)\}
  +
  \{\E(\hat\theta)-\theta\}.
\]
因此
\[
\begin{aligned}
  \MSE(\hat\theta)
  &=
  \E[(\hat\theta-\theta)^2]\\
  &=
  \E[(\hat\theta-\E\hat\theta)^2]
  +2b\,\E[\hat\theta-\E(\hat\theta)]
  +b^2.
\end{aligned}
\]
由于
\[
  \E[\hat\theta-\E(\hat\theta)]=0,
\]
交叉项为零，故
\[
  \boxed{
  \MSE(\hat\theta)
  =
  \Var(\hat\theta)+\Bias(\hat\theta)^2.
  }
\]

\textbf{（2）$T$ 的 bias 与 MSE。}

\[
\begin{aligned}
  \Bias(T)
  &=
  \E(T)-\theta\\
  &=
  2\theta+1-\theta\\
  &=
  \boxed{\theta+1}.
\end{aligned}
\]
所以
\[
  \boxed{
  \MSE(T)
  =
  \frac9n+(\theta+1)^2.
  }
\]

\textbf{（3）构造无偏估计量。}

由
\[
  \E(T)=2\theta+1,
\]
自然消去常数项并除以系数 2：
\[
  \boxed{U=\frac{T-1}{2}}.
\]
验证：
\[
  \E(U)
  =
  \frac{\E(T)-1}{2}
  =
  \frac{2\theta+1-1}{2}
  =
  \theta.
\]
所以 $U$ 无偏。方差为
\[
  \Var(U)
  =
  \frac14\Var(T)
  =
  \boxed{\frac9{4n}}.
\]
因为 $U$ 无偏，
\[
  \boxed{\MSE(U)=\frac9{4n}}.
\]

\textbf{（4）组合两个无偏估计量。}

因为 $U,S$ 都无偏，
\[
  \E(W_c)=c\theta+(1-c)\theta=\theta,
\]
故 $\MSE(W_c)=\Var(W_c)$。利用独立性：
\[
\begin{aligned}
  \Var(W_c)
  &=
  c^2\frac9{4n}
  +(1-c)^2\frac3n.
\end{aligned}
\]
去掉公共因子 $1/n$，对
\[
  g(c)=\frac94c^2+3(1-c)^2
\]
求导：
\[
  g'(c)=\frac92c-6(1-c)=\frac{21}{2}c-6.
\]
令 $g'(c)=0$：
\[
  \boxed{c^\star=\frac{12}{21}=\frac47}.
\]
二阶导
\[
  g''(c)=\frac{21}{2}>0,
\]
所以这是最小值。代回：
\[
\begin{aligned}
  \MSE(W_{c^\star})
  &=
  \frac9{4n}\left(\frac47\right)^2
  \frac3n\left(\frac37\right)^2\\
  &=
  \frac{36}{49n}+\frac{27}{49n}\\
  &=
  \boxed{\frac9{7n}}.
\end{aligned}
\]
最小 MSE 小于 $9/(4n)$ 和 $3/n$，因此合适的组合比单独使用任一估计量都更精确。
\end{solutionbox}

\begin{solutionbox}[第 2 题解析：单总体均值与总体比例]
\textbf{第一部分：总体均值。}

总体正态、方差未知，使用单样本 $t$ 推断：
\[
  \nu=n-1=24,
  \qquad
  \SE(\bar X)=\frac{s}{\sqrt n}=\frac5{\sqrt{25}}=1.
\]

\textbf{（1）95\% 置信区间。}

\[
\begin{aligned}
  \bar x
  \pm
  t_{0.025,24}\frac{s}{\sqrt n}
  &=
  41.2\pm2.064(1)\\
  &=
  \boxed{(39.136,\ 43.264)}.
\end{aligned}
\]
解释：我们有 95\% 的置信度认为总体均值 $\mu$ 位于 39.136 到 43.264 之间。更严格地说，若重复抽样并用同一规则构造区间，约 95\% 的区间会覆盖真实均值。

\textbf{（2）右尾检验。}

原假设边界为 $\mu_0=40$：
\[
  t_{\mathrm{obs}}
  =
  \frac{\bar x-\mu_0}{s/\sqrt n}
  =
  \frac{41.2-40}{1}
  =
  1.2.
\]
在 $H_0$ 边界及正态总体假设下，
\[
  T\sim t_{24}.
\]
右尾 $\alpha=0.05$ 的拒绝域为
\[
  t_{\mathrm{obs}}>t_{0.05,24}=1.711.
\]
因为
\[
  1.2<1.711,
\]
所以\textbf{不能拒绝 $H_0$}。样本没有提供足够证据说明总体均值大于 40。

\textbf{第二部分：总体比例。}

\[
  \hat p=\frac{116}{400}=0.29.
\]

\textbf{（3）95\% Wald 置信区间。}

区间标准误使用样本比例 $\hat p$：
\[
  \SE_{\mathrm{CI}}
  =
  \sqrt{\frac{\hat p(1-\hat p)}n}
  =
  \sqrt{\frac{0.29(0.71)}{400}}
  \approx0.022688.
\]
因此
\[
\begin{aligned}
  \hat p\pm1.96\SE_{\mathrm{CI}}
  &=
  0.29\pm1.96(0.022688)\\
  &=
  \boxed{(0.2455,\ 0.3345)}.
\end{aligned}
\]

\textbf{（4）双尾比例检验。}

检验时必须在 $H_0:p=p_0=0.25$ 下计算标准误：
\[
  \SE_0
  =
  \sqrt{\frac{0.25(0.75)}{400}}
  \approx0.021651.
\]
\[
  Z_{\mathrm{obs}}
  =
  \frac{0.29-0.25}{0.021651}
  \approx1.848.
\]
双尾 $\alpha=0.05$ 的拒绝域为
\[
  |Z_{\mathrm{obs}}|>1.960.
\]
因为 $1.848<1.960$，不能拒绝 $H_0$。没有足够证据说明总体比例与 0.25 不同。

区间和检验使用不同比例的原因是：区间中 $p$ 未被指定，只能用 $\hat p$ 估计标准误；检验中假定 $H_0$ 为真，因此抽样方差应使用原假设值 $p_0$。
\end{solutionbox}

\begin{solutionbox}[第 3 题解析：配对样本与两个独立大样本]
\textbf{第一部分：配对样本。}

\textbf{（1）为什么使用 paired-sample 方法。}

改革前后成绩来自同一个学生，因此每个“改革后成绩”都与该学生自己的“改革前成绩”天然配对。两个观测通常相关，不能按两个独立样本处理。应先构造
\[
  D_i=(\text{改革后})-(\text{改革前}),
\]
然后对差值总体均值 $\mu_D$ 进行单样本 $t$ 推断。

\textbf{（2）右尾配对 $t$ 检验。}

\[
  \SE(\bar D)
  =
  \frac{s_D}{\sqrt n}
  =
  \frac3{\sqrt{16}}
  =
  0.75.
\]
\[
  t_{\mathrm{obs}}
  =
  \frac{\bar d-0}{s_D/\sqrt n}
  =
  \frac{2.4}{0.75}
  =
  3.2.
\]
自由度为
\[
  \nu=n-1=15.
\]
右尾拒绝域：
\[
  t_{\mathrm{obs}}>t_{0.05,15}=1.753.
\]
因为 $3.2>1.753$，拒绝 $H_0$。样本提供了足够证据说明教学改革提高了平均成绩。

\textbf{（3）$\mu_D$ 的 95\% 置信区间。}

\[
\begin{aligned}
  \bar d
  \pm
  t_{0.025,15}\frac{s_D}{\sqrt n}
  &=
  2.4\pm2.131(0.75)\\
  &=
  \boxed{(0.802,\ 3.998)}.
\end{aligned}
\]
区间整体为正，表示平均提高幅度估计在约 0.802 到 3.998 分之间。

\textbf{第二部分：两个独立大样本。}

两个样本独立且样本量较大，使用大样本正态近似：
\[
\begin{aligned}
  \SE(\bar X_1-\bar X_2)
  &=
  \sqrt{\frac{4^2}{100}+\frac{3.6^2}{81}}\\
  &=
  \sqrt{0.16+0.16}\\
  &=
  \sqrt{0.32}
  \approx0.5657.
\end{aligned}
\]
点估计：
\[
  \bar x_1-\bar x_2=12.8-11.5=1.3.
\]

\textbf{（4）95\% 置信区间。}

\[
\begin{aligned}
  1.3\pm1.96(0.5657)
  &=
  \boxed{(0.191,\ 2.409)}.
\end{aligned}
\]

\textbf{（5）双尾检验。}

\[
  Z_{\mathrm{obs}}
  =
  \frac{1.3-0}{0.5657}
  \approx2.298.
\]
拒绝域：
\[
  |Z_{\mathrm{obs}}|>1.960.
\]
因为 $2.298>1.960$，拒绝 $H_0$。样本提供了足够证据说明两个总体均值不同。由于点估计和整个置信区间均为正，差异方向为 $\mu_1>\mu_2$。
\end{solutionbox}

\begin{solutionbox}[第 4 题解析：简单回归计算与多元回归证明]
\textbf{第一部分：简单线性回归。}

\textbf{（1）拟合、平方和与决定系数。}

\[
  \hat\beta_1
  =
  \frac{S_{xy}}{S_{xx}}
  =
  \frac{200}{80}
  =
  2.5.
\]
\[
  \hat\beta_0
  =
  \bar y-\hat\beta_1\bar x
  =
  30-2.5(12)
  =
  0.
\]
所以拟合直线为
\[
  \boxed{\hat y=2.5x}.
\]
回归平方和：
\[
  SSR
  =
  \frac{S_{xy}^2}{S_{xx}}
  =
  \frac{200^2}{80}
  =
  500.
\]
\[
  SSE
  =
  S_{yy}-SSR
  =
  620-500
  =
  \boxed{120}.
\]
\[
  s_e^2
  =
  \frac{SSE}{n-2}
  =
  \frac{120}{8}
  =
  \boxed{15}.
\]
\[
  R^2
  =
  \frac{SSR}{S_{yy}}
  =
  \frac{500}{620}
  \approx
  \boxed{0.806}.
\]
约 80.6\% 的响应总变异由该简单线性回归模型解释。

\textbf{（2）斜率检验。}

\[
  \SE(\hat\beta_1)
  =
  \frac{s_e}{\sqrt{S_{xx}}}
  =
  \sqrt{\frac{15}{80}}
  \approx0.4330.
\]
\[
  t_{\mathrm{obs}}
  =
  \frac{2.5-0}{0.4330}
  \approx5.774.
\]
自由度为 $n-2=8$。双尾拒绝域：
\[
  |t_{\mathrm{obs}}|>t_{0.025,8}=2.306.
\]
因为 $5.774>2.306$，拒绝 $H_0$；数据支持总体线性斜率不为零。

\textbf{（3）平均响应 CI 与个体预测 PI。}

在 $x_0=14$：
\[
  \hat y_0
  =
  2.5(14)
  =
  35.
\]
平均响应的估计标准误：
\[
\begin{aligned}
  \SE_{\mathrm{mean}}
  &=
  s_e\sqrt{\frac1n+\frac{(x_0-\bar x)^2}{S_{xx}}}\\
  &=
  \sqrt{15}
  \sqrt{\frac1{10}+\frac{(14-12)^2}{80}}\\
  &=
  \sqrt{15(0.15)}
  =
  1.5.
\end{aligned}
\]
所以 95\% mean-response CI 为
\[
  35\pm2.306(1.5)
  =
  \boxed{(31.541,\ 38.459)}.
\]

新个体预测的估计标准误：
\[
\begin{aligned}
  \SE_{\mathrm{pred}}
  &=
  s_e\sqrt{
    1+\frac1n+\frac{(x_0-\bar x)^2}{S_{xx}}
  }\\
  &=
  \sqrt{15(1.15)}
  \approx4.153.
\end{aligned}
\]
所以 95\% prediction interval 为
\[
  35\pm2.306(4.153)
  =
  \boxed{(25.422,\ 44.578)}.
\]
prediction interval 更宽，因为预测一个新个体时，除了回归均值估计的不确定性，还包含该新个体自身的随机误差；公式中因此多出一个 $1$。

\textbf{第二部分：多元线性回归理论。}

\textbf{（4）推导矩阵最小二乘与残差正交。}

展开目标函数：
\[
\begin{aligned}
  Q(\mathbf b)
  &=
  (\by-\bX\mathbf b)^\mathsf{T}(\by-\bX\mathbf b)\\
  &=
  \by^\mathsf{T}\by
  -2\mathbf b^\mathsf{T}\bX^\mathsf{T}\by
  +\mathbf b^\mathsf{T}\bX^\mathsf{T}\bX\mathbf b.
\end{aligned}
\]
对 $\mathbf b$ 求梯度：
\[
  \nabla_{\mathbf b}Q
  =
  -2\bX^\mathsf{T}\by
  +2\bX^\mathsf{T}\bX\mathbf b.
\]
令梯度为零，得到正规方程：
\[
  \bX^\mathsf{T}\bX\hat{\bbeta}
  =
  \bX^\mathsf{T}\by.
\]
由于 $\bX$ 满列秩，$\bX^\mathsf{T}\bX$ 可逆：
\[
  \boxed{
  \hat{\bbeta}
  =
  (\bX^\mathsf{T}\bX)^{-1}\bX^\mathsf{T}\by.
  }
\]
令
\[
  \be=\by-\bX\hat{\bbeta}.
\]
由正规方程：
\[
  \bX^\mathsf{T}\be
  =
  \bX^\mathsf{T}\by-\bX^\mathsf{T}\bX\hat{\bbeta}
  =
  \boxed{\mathbf0}.
\]
即残差向量与设计矩阵每一列都正交。

\textbf{（5）无偏性与方差矩阵。}

代入 $\by=\bX\bbeta+\bvarepsilon$：
\[
\begin{aligned}
  \hat{\bbeta}
  &=
  (\bX^\mathsf{T}\bX)^{-1}
  \bX^\mathsf{T}
  (\bX\bbeta+\bvarepsilon)\\
  &=
  \bbeta
  +
  (\bX^\mathsf{T}\bX)^{-1}
  \bX^\mathsf{T}\bvarepsilon.
\end{aligned}
\]
取期望：
\[
  \E(\hat{\bbeta})
  =
  \bbeta+
  (\bX^\mathsf{T}\bX)^{-1}
  \bX^\mathsf{T}\E(\bvarepsilon)
  =
  \boxed{\bbeta}.
\]
令
\[
  \mathbf A=(\bX^\mathsf{T}\bX)^{-1}\bX^\mathsf{T}.
\]
则
\[
\begin{aligned}
  \Var(\hat{\bbeta})
  &=
  \mathbf A\Var(\bvarepsilon)\mathbf A^\mathsf{T}\\
  &=
  \sigma^2
  (\bX^\mathsf{T}\bX)^{-1}
  \bX^\mathsf{T}\bX
  (\bX^\mathsf{T}\bX)^{-1}\\
  &=
  \boxed{
  \sigma^2(\bX^\mathsf{T}\bX)^{-1}.
  }
\end{aligned}
\]
\end{solutionbox}

\begin{solutionbox}[第 5 题解析：四类分布的 MLE 与 Fisher 信息]
\textbf{（1）Normal：$\mu$ 与 $\sigma^2$ 都未知。}

令 $v=\sigma^2$。忽略与参数无关的常数，log-likelihood 为
\[
  \ell(\mu,v)
  =
  -\frac n2\log v
  -
  \frac1{2v}\sum_{i=1}^n(y_i-\mu)^2
  +\text{constant}.
\]
score components 为
\[
  \frac{\partial\ell}{\partial\mu}
  =
  \frac1v\sum_{i=1}^n(y_i-\mu),
\]
\[
  \frac{\partial\ell}{\partial v}
  =
  -\frac n{2v}
  +
  \frac1{2v^2}\sum_{i=1}^n(y_i-\mu)^2.
\]
解 score equations：
\[
  \boxed{\hat\mu_{\mathrm{MLE}}=\bar y=5},
\]
\[
  \boxed{
  \hat\sigma^2_{\mathrm{MLE}}
  =
  \frac1n\sum_{i=1}^n(y_i-\bar y)^2
  =
  \frac{24}{8}
  =
  3.
  }
\]
方差 MLE 的分母为 $n$，因为这是最大化 likelihood 得到的解。无偏样本方差使用 $n-1$，是为了修正估计均值后损失的一个自由度；二者优化标准不同。

\textbf{（2）Binomial：边界 MLE。}

当 $m=20,y=0$：
\[
  L(p)
  =
  \binom{20}{0}p^0(1-p)^{20}
  =
  (1-p)^{20},
  \qquad
  0\le p\le1.
\]
在 $0<p<1$ 上，
\[
  \ell'(p)
  =
  -\frac{20}{1-p}<0.
\]
所以 likelihood 随 $p$ 增大而严格下降，最大值出现在左端点：
\[
  \boxed{\hat p_{\mathrm{MLE}}=0}.
\]
本题 score equation 在参数空间内部没有根，因此若只机械地令 score 为零会漏掉答案，必须检查边界。

形式上代入 Wald 标准误公式会得到
\[
  \sqrt{\frac{\hat p(1-\hat p)}{20}}=0,
\]
但这并不代表估计没有不确定性。参数位于边界时，MLE 的普通内部渐近正态近似失效，因此普通 Wald 区间或检验不可靠。

\textbf{（3）Poisson MLE 与 Wald 区间。}

\[
  \ell(\lambda)
  =
  -n\lambda
  +\left(\sum y_i\right)\log\lambda
  +\text{constant}.
\]
\[
  U(\lambda)
  =
  -n+\frac{\sum y_i}{\lambda}.
\]
令 $U(\lambda)=0$：
\[
  \boxed{
  \hat\lambda
  =
  \frac{\sum y_i}{n}
  =
  \frac{64}{16}
  =
  4.
  }
\]
样本 Fisher 信息：
\[
  \boxed{I_n(\lambda)=\frac n\lambda}.
\]
代入 MLE：
\[
  I_n(\hat\lambda)=\frac{16}{4}=4,
\]
\[
  \widehat{\SE}(\hat\lambda)
  =
  \frac1{\sqrt{I_n(\hat\lambda)}}
  =
  \boxed{0.5}.
\]
95\% Wald CI：
\[
  4\pm1.96(0.5)
  =
  \boxed{(3.020,\ 4.980)}.
\]

\textbf{（4）Exponential rate MLE 与 Wald 检验。}

密度为
\[
  f(y;\lambda)=\lambda e^{-\lambda y},
  \qquad y\ge0,\quad\lambda>0.
\]
log-likelihood：
\[
  \ell(\lambda)
  =
  n\log\lambda-\lambda\sum y_i.
\]
score：
\[
  U(\lambda)
  =
  \frac n\lambda-\sum y_i.
\]
所以
\[
  \boxed{
  \hat\lambda
  =
  \frac n{\sum y_i}
  =
  \frac{16}{40}
  =
  0.4.
  }
\]
样本 Fisher 信息：
\[
  \boxed{
  I_n(\lambda)=\frac n{\lambda^2}.
  }
\]
plug-in 标准误：
\[
  \widehat{\SE}(\hat\lambda)
  =
  \frac{\hat\lambda}{\sqrt n}
  =
  \frac{0.4}{4}
  =
  0.1.
\]
Wald 统计量：
\[
  Z_{\mathrm{obs}}
  =
  \frac{\hat\lambda-0.6}
  {\widehat{\SE}(\hat\lambda)}
  =
  \frac{0.4-0.6}{0.1}
  =
  -2.
\]
双尾 $\alpha=0.05$ 的拒绝域为
\[
  |Z_{\mathrm{obs}}|>1.960.
\]
因为
\[
  |-2|>1.960,
\]
所以拒绝 $H_0$。在该 plug-in Wald 近似下，数据提供了足够证据说明 exponential rate 与 0.6 不同。
\end{solutionbox}

\end{document}
